\section{Discussion}
\label{sec:discussion}

\subsection{What the Live Salmon Test Reveals}
\label{sec:interpretation}

Our results paint a nuanced picture of \aiscientist capabilities.
On one hand, these systems are statistically conservative and competent at standard analysis: they reliably detect strong linear effects, never hallucinate significance on null data, and consistently apply multiple comparisons corrections.
On the other hand, they exhibit a systematic blind spot for nonlinear structure, applying the same correlation-to-regression pipeline regardless of whether it is sufficient for the data at hand.

The 0\% interaction detection rate is particularly informative because it is not a failure of statistical power.
The $X_4 \times X_5$ interaction has a coefficient of 0.8 and contributes 8\% of total variance---a meaningful signal that any analyst who tested for interactions would find.
The failure is one of \textit{methodology}, not \textit{sensitivity}: \aiscientists do not attempt the analytical steps that would reveal nonlinear relationships.
This suggests that current \llms lack what we might call ``analytical creativity''---the ability to move beyond default workflows when data characteristics warrant it.

\subsection{Implications for AI-Generated Research}
\label{sec:implications}

\para{For developers of \aiscientist systems.}
The bottleneck we identify is not statistical rigor (which is adequate) but analytical breadth.
Future systems should incorporate systematic model diagnostic steps: residual analysis after initial linear fits, automated interaction screening, and comparison of linear vs.\ nonlinear models.
These are standard practices in competent data analysis but are absent from the \llm ``default workflow.''

\para{For consumers of AI-generated research.}
Our results suggest that AI-generated papers should be treated as first-pass linear analyses.
Nonlinear effects, interaction terms, and complex functional forms may be systematically underreported.
When the collective \aiscientist literature on a dataset uniformly agrees on a linear model, this consensus may reflect shared methodological limitations rather than genuine evidence for linearity.

\para{For the research community.}
The live salmon test methodology---providing known-\dgp data and evaluating recovery---is a valuable complement to existing benchmarks.
While benchmarks like PaperBench \citep{starace2025paperbench} and BLADE \citep{gu2024blade} measure broad research capability, the live salmon test provides a precise, controlled measurement of what specific patterns \aiscientists can and cannot detect.
We encourage the development of a diverse suite of live salmon tests covering different data types, effect structures, and analytical challenges.

\subsection{Connection to the Dead Salmon Tradition}
\label{sec:tradition}

Our dead salmon results are the inverse of \citet{bennett2009neural}: where standard fMRI analysis produced false positives on null data, \aiscientists produce zero false positives.
This difference likely arises because \llms have internalized the lesson of the original dead salmon study---multiple comparisons correction is well-represented in statistical training data and is a ``textbook'' step that \llms reliably reproduce.
The broader implication is that \llms are well-calibrated for \textit{known} failure modes (false positives from multiple testing) but poorly equipped for \textit{novel} analytical challenges (detecting unexpected nonlinear structure).

This finding parallels \citet{meloux2025dead}, who showed that AI interpretability tools produce plausible explanations for random networks.
Both results highlight that AI analytical pipelines can be simultaneously conservative in some respects and blind in others---a pattern that is harder to diagnose than simple unreliability.

\subsection{Limitations}
\label{sec:limitations}

Our study has several limitations that should be considered when interpreting the results.

\para{Data delivery format.}
We provide summary statistics rather than raw data, and models cannot execute code.
Full \aiscientist systems like The AI Scientist \citep{lu2024ai} can write and run analysis code, which might enable more diverse methods.
However, the summary statistics we provide are sufficient for a competent analyst to suspect nonlinear structure (\eg by noting that $R^2$ is well below 1.0 after accounting for linear terms).

\para{Prompt design.}
Our prompt asks for general analysis without specifically mentioning interactions.
A more directive prompt (\eg ``check for interaction effects'') might change results.
This is by design: we test whether \aiscientists independently identify the need for interaction screening, which is the capability gap we aim to measure.

\para{Single-shot paradigm.}
Real \aiscientist systems iterate: they generate hypotheses, run code, examine results, and revise.
Our single-shot evaluation may underestimate the capability of iterative systems.
However, single-shot analysis tests the core analytical reasoning that drives the first iteration of any multi-step system.

\para{Evaluation by \llm.}
We use \gptmodel to extract claims from generated papers, which may introduce parsing errors.
We mitigate this by using temperature 0 and structured extraction prompts, but human verification of a subset would strengthen confidence.

\para{Limited scale.}
With 3 models and 3 runs each, we have limited statistical power for cross-model comparisons.
The core finding (0/9 interaction detection) is robust due to its extreme value, but subtler differences between models require larger samples to assess reliably.

\para{Single \dgp.}
We test a single data generating process with specific effect sizes and feature correlations.
The interaction detection failure may vary with different signal-to-noise ratios, interaction types, or data dimensionalities.
A comprehensive benchmark would include a suite of \dgp{}s spanning diverse analytical challenges.
